\begin{center}
\SourceLeaf{027b}\pdfbookmark[2]{9.}{article-009}9.
\end{center}
Si perpendiculum e puncto quocunque ellipsis in lineam apsidum demissum retro producitur, usquedum circulo e centro ellipsis radio $a$ descripto occurrat, inclinatio eius radii, qui puncto intersectionis respondet, contra lineam apsidum (simili modo intellecta ut supra pro anomalia vera) anomaliae excentricae aequalis erit, ut nullo negotio ex aequ. IX art. praec. deducitur. Porro patet, $r\sin v$ esse distantiam cuiusque puncti ellipsis a linea apsidum; quae quum per aequ. VIII fiat $= a\cos\varphi\sin E$, maxima erit pro $E = 90^0$, i. e. in centro ellipsis. Haecce distantia maxima, quae fit $= a\cos\varphi = \frac{p}{\cos\varphi} = \BareRoot{ap}$, \emph{semiaxis minor} appellatur. In foco ellipsis, i. e. pro $v = 90^0$, distantia ista manifesto fit $= p$, sive semiparametro aequalis.
